\documentclass[11pt,a4paper]{article}

% ---------- Encoding & Sprache ----------
\usepackage[utf8]{inputenc}
\usepackage[T1]{fontenc}
\usepackage[ngerman]{babel}
\usepackage{lmodern}

% ---------- Layout ----------
\usepackage[margin=2cm]{geometry}
\usepackage{setspace}
\usepackage{enumitem}
\setlist{nosep, itemsep=1pt, topsep=2pt}

\setlength{\parindent}{0pt}
\setlength{\parskip}{4pt}

\title{\vspace{-1.5cm}\Large\textbf{Portfolioteil 2 -- Erarbeitungs-/Reflexionsphase}\\[2pt]
\large Kurze Zusammenfassung der Implementierung\\[2pt]
\normalsize DLBDSPBDM01\_D -- Data-Mart-Erstellung in SQL}
\author{Tizian Senger \\ Matrikelnummer: IU14143428}
\date{\today}

\begin{document}
\maketitle
\vspace{-1cm}

% =========================================================
\section*{1. Verwendetes Datenbankmanagementsystem}
% =========================================================
% TODO: Welches DBMS wurde gewaehlt (z. B. MySQL/MariaDB/PostgreSQL) und warum?

% =========================================================
\section*{2. Umsetzung des ER-Modells}
% =========================================================
% TODO: Kurz beschreiben, wie die 11 Entitaeten aus Phase 1 als Tabellen umgesetzt wurden,
% wie die binaeren M:N-Beziehungen (Verfasst, Gehoert zu) als Verknuepfungstabellen aufgeloest wurden
% und wie die drei Dreifachbeziehungen (Leiht_aus, Bewertet, Verfuegbar_an) als eigene Tabellen
% mit Fremdschluesseln auf alle drei beteiligten Tabellen realisiert wurden.

% =========================================================
\section*{3. Normalisierung}
% =========================================================
% TODO: Kurz begruenden, warum das Schema in der 3. Normalform vorliegt
% (keine transitiven Abhaengigkeiten, Auslagerung von Autor/Verlag/Genre/Sprache in eigene Tabellen).

% =========================================================
\section*{4. Dummy-Daten und Testfaelle}
% =========================================================
% TODO: Kurz beschreiben, wie die Testdaten erzeugt wurden (mind. 10 Eintraege je Tabelle)
% und welchen Testfall/welche Abfrage(n) zur Ueberpruefung der Datenbank verwendet wurden.

% =========================================================
\section*{5. Herausforderungen und Erkenntnisse}
% =========================================================
% TODO: Kurze Reflexion: Was war beim Implementieren schwierig? Was wuerde man anders machen?

\end{document}
